% Generated draft: requires mathematical and source review.
\documentclass[11pt]{article}
\usepackage[margin=25mm]{geometry}
\usepackage{fontspec}
\setmainfont{Latin Modern Roman}
\usepackage{amsmath,amssymb,mathtools,unicode-math}
\setmathfont{Latin Modern Math}
\usepackage{enumitem,needspace,xurl,hyperref}
\setlength{\parindent}{0pt}
\setlength{\parskip}{5pt}
\setlength{\emergencystretch}{3em}
\newcommand{\N}{\mathbb{N}}
\newcommand{\Z}{\mathbb{Z}}
\newcommand{\Q}{\mathbb{Q}}
\newcommand{\R}{\mathbb{R}}
\newcommand{\C}{\mathbb{C}}
\newcommand{\F}{\mathbb{F}}
\newcommand{\A}{\mathbb{A}}
\newcommand{\PP}{\mathbb{P}}
\newcommand{\E}{\mathbb{E}}
\newcommand{\eps}{\varepsilon}
\newcommand{\dd}{\,\mathrm{d}}
\begin{document}
\section*{UnsolvedMath ID: 800045}
\textbf{Source problem code:} catalogue entry
\subsection{Definitions and mathematical statement}
A semigroup is a nonempty set equipped with an associative binary operation. A semigroup variety is a class of semigroups defined by identities in this binary operation. For a semigroup \(S\), let \(\operatorname{var}(S)\) denote the smallest semigroup variety containing \(S\). A subvariety of a variety \(\mathcal V\) is a semigroup variety \(\mathcal W\subseteq\mathcal V\).

The supplied record does not define the term \emph{inclusive variety}. Consequently, the following is a conditional formalization with that term left unresolved: let \(\mathfrak I\) denote the class of semigroup varieties called inclusive in the source, and put
\[
\mathcal L_{\mathfrak I}(\mathcal V)
=\{\mathcal W:\mathcal W\text{ is a semigroup variety},\ \mathcal W\subseteq\mathcal V,\ \mathcal W\in\mathfrak I\},
\]
ordered by inclusion. The source calls this ordered collection the lattice of inclusive subvarieties; its lattice structure cannot be specified further without the missing definition.

\textbf{Question (G. Mashevitzky; source Problem 47).} Does there exist a finite semigroup \(S\) such that
\[
\operatorname{var}(S)\in\mathfrak I
\qquad\text{and}\qquad
\mathcal L_{\mathfrak I}(\operatorname{var}(S))\text{ is infinite}?
\]
The intended class \(\mathfrak I\) must be identified before this becomes a fully specified mathematical question.

The supplied background asserts that if an inclusive variety is generated by a finite group or by a finite nilsemigroup, then its lattice of inclusive subvarieties is finite, citing \([\mathrm{M3}]\). Here a group is regarded as a semigroup under its multiplication, and a nilsemigroup means a semigroup with an absorbing zero \(0\) such that for every element \(x\) there is an integer \(n\geq 1\) with \(x^n=0\). This background assertion is reported from the source and has not been independently verified; the bibliographic identity of \([\mathrm{M3}]\) is not supplied.
\subsection{Short English statement}
G. Mashevitzky asks whether a finite semigroup can generate an inclusive variety having infinitely many inclusive subvarieties. The record reports finiteness for inclusive varieties generated by finite groups or finite nilsemigroups. A complete formulation requires the missing definition of \emph{inclusive variety}. The registry labels the problem open; that status is supplied metadata.
\subsection{Sources}
\begin{itemize}
\item Local UnsolvedMath registry, permanent ID \(800045\), exact source code \texttt{catalogue entry}, title \emph{The Lyapin's Notebook --- Problem 47}. The supplied record labels the problem open and attributes the question itself to G. Mashevitzky.
\item Bershadsky, Kublanovsky, and Mashevitzky, \emph{The Lyapin's Notebook: Unsolved Problems in Semigroup Theory} (2026), Problem 47. Supplied source URL: \url{https://arxiv.org/abs/2604.04763}. The record reports access on 2026-10-16 and extraction by \texttt{arxiv\_latex}; these are supplied provenance details, not an independent source visit or verification.
\item Reference \([\mathrm{M3}]\), cited in the supplied statement for the finiteness assertion concerning finite groups and finite nilsemigroups. No bibliographic details or link for this reference are supplied.
\end{itemize}
\end{document}
